\section{Código Modelica de los subsistemas añadidos}
\label{sec:codigo}

Se muestran las ecuaciones de cada componente. Las anotaciones gráficas (iconos y
posición de los conectores) se omiten por brevedad.

\subsection{\texttt{SubSistemas.ModeloMotor}}
\begin{lstlisting}
model ModeloMotor "Equivalent DC motor + propeller load (one rotor)"
  parameter SI.Resistance R "Armature resistance";
  parameter SI.Inductance L "Armature inductance";
  parameter Real ke(unit = "V.s/rad") "Back-EMF constant";
  parameter Real kt(unit = "N.m/A") = ke "Torque constant (= ke in SI units)";
  parameter SI.Inertia Jm "Rotor inertia (motor + propeller)";
  parameter Real bm(unit = "N.m.s/rad") "Viscous friction coefficient";
  parameter Real km(unit = "N.m.s2/rad2") "Propeller torque coefficient";
  parameter Boolean steadyStateInit = true;
  parameter SI.AngularVelocity w_start = 0;

  RealInput  V(unit = "V") "Applied (equivalent) voltage";
  RealOutput w(unit = "rad/s", start = w_start) "Rotor angular speed";
  RealOutput i(unit = "A", start = km*w_start^2/kt) "Armature current";

  SI.Torque tau_m "Electromagnetic torque";
  SI.Torque tau_p "Propeller drag torque";
  SI.Power  P_el  "Electrical power drawn";
initial equation
  if steadyStateInit then
    der(i) = 0;
    der(w) = 0;
  else
    i = 0;
    w = w_start;
  end if;
equation
  L*der(i) = V - R*i - ke*w;            // modelo eléctrico
  tau_m = kt*i;
  tau_p = km*w*abs(w);                  // = km w^2 para w >= 0
  Jm*der(w) = tau_m - tau_p - bm*w;     // modelo mecánico
  P_el = V*i;
end ModeloMotor;
\end{lstlisting}

\subsection{\texttt{SubSistemas.ESC}}
\begin{lstlisting}
block ESC "Speed controller: rotor speed command -> motor voltage (4 rotors)"
  parameter SI.Resistance R;
  parameter Real ke(unit = "V.s/rad");
  parameter Real kt(unit = "N.m/A") = ke;
  parameter Real bm(unit = "N.m.s/rad");
  parameter Real km(unit = "N.m.s2/rad2");
  parameter Real kp(unit = "V.s/rad") = 0.05 "Proportional gain";
  parameter Real ki(unit = "V/rad")   = 0.75 "Integral gain";
  parameter Boolean steadyStateInit = true;

  RealInput  w_ref[4](each unit = "rad/s") "Rotor speed commands";
  RealInput  w_m[4](each unit = "rad/s")   "Measured rotor speeds";
  RealInput  V_bat(unit = "V")             "Available battery voltage";
  RealOutput V[4](each unit = "V")         "Motor voltages, 0 <= V <= V_bat";

  Real e[4], x[4](each start = 0);
  SI.Voltage V_ff[4], V_u[4];
initial equation
  if steadyStateInit then der(x) = zeros(4); else x = zeros(4); end if;
equation
  for j in 1:4 loop
    e[j]    = w_ref[j] - w_m[j];
    V_ff[j] = ke*w_ref[j] + R/kt*(km*w_ref[j]^2 + bm*w_ref[j]);
    V_u[j]  = V_ff[j] + kp*e[j] + ki*x[j];
    V[j]    = min(max(V_u[j], 0), V_bat);
    der(x[j]) = if (V_u[j] > V_bat and e[j] > 0) or (V_u[j] < 0 and e[j] < 0)
                then 0 else e[j];          // anti-windup
  end for;
end ESC;
\end{lstlisting}

\subsection{\texttt{SubSistemas.Bateria}}
\begin{lstlisting}
model Bateria "Battery as an energy reservoir with constant voltage until empty"
  parameter SI.Voltage V_nom = 7.4 "Nominal voltage";
  parameter Real C_Ah(unit = "A.h") = 0.5 "Nominal capacity";
  parameter SI.Energy E_0 = V_nom*C_Ah*3600 "Initial stored energy";

  RealInput  V[4](each unit = "V") "Motor voltages";
  RealInput  i[4](each unit = "A") "Motor currents";
  RealOutput V_bat(unit = "V")     "Available voltage";
  RealOutput SOC(unit = "1")       "State of charge E/E_0";

  SI.Energy E(start = E_0, fixed = true) "Stored energy";
  SI.Power  P "Electrical power drawn by the motors";
  Boolean empty = E <= 0;
equation
  P = sum(V[j]*i[j] for j in 1:4);
  der(E) = -P;
  V_bat = if empty then 0 else V_nom;
  SOC = E/E_0;
end Bateria;
\end{lstlisting}

\subsection{\texttt{SubSistemas.Perturbaciones}}
\begin{lstlisting}
model Perturbaciones "External forces in the inertial frame: drag + wind"
  parameter Real c_d(unit = "N.s/m") = 0 "Linear drag coefficient";
  parameter SI.Force F_wind[3] = {0, 0, 0} "Wind force, world frame";
  parameter SI.Time t_wind = 0 "Wind start time";

  Modelica.Mechanics.MultiBody.Interfaces.Frame_b frame_b;

  SI.Velocity v[3]; SI.Force F_d[3], F_w[3], F[3];
equation
  v   = der(frame_b.r_0);
  F_d = -c_d*v;
  F_w = if time >= t_wind then F_wind else zeros(3);
  F   = F_d + F_w;
  frame_b.f = -Frames.resolve2(frame_b.R, F);   // como WorldForce
  frame_b.t = zeros(3);
end Perturbaciones;
\end{lstlisting}

\subsection{Cambios en \texttt{Drone}}
\begin{lstlisting}
parameter SI.Mass m_payload = 0 "Carried payload";
parameter SI.Resistance R_mot = 0.1;  parameter SI.Inductance L_mot = 1e-4;
parameter Real ke(unit = "V.s/rad") = 0.005;
parameter SI.Inertia Jm = 2e-5;  parameter Real bm(unit = "N.m.s/rad") = 1e-6;
parameter SI.Voltage V_bat = 7.4;  parameter Real C_bat(unit = "A.h") = 0.5;
parameter Real c_d(unit = "N.s/m") = 0;
parameter SI.Force F_wind[3] = {0, 0, 0};  parameter SI.Time t_wind = 0;
parameter SI.AngularVelocity w_hover = sqrt((mass + m_payload)*g_n/(4*kf));

SubSistemas.ESC esc(R = R_mot, ke = ke, bm = bm, km = km);
SubSistemas.ModeloMotor motor[4](each R = R_mot, each L = L_mot, each ke = ke,
  each Jm = Jm, each bm = bm, each km = km, each w_start = w_hover);
SubSistemas.Bateria bateria(V_nom = V_bat, C_Ah = C_bat);
SubSistemas.Perturbaciones perturbaciones(c_d = c_d, F_wind = F_wind, t_wind = t_wind);
Modelica.Mechanics.MultiBody.Parts.Body body(m = mass + m_payload, ...,
  r_0(each fixed = true), v_0(each fixed = true),
  angles_fixed = true, w_0_fixed = true);
\end{lstlisting}
